\documentclass[conference]{IEEEtran}
\IEEEoverridecommandlockouts

% Paquetes necesarios
\usepackage{cite}
\usepackage{amsmath,amssymb,amsfonts}
\usepackage{algorithmic}
\usepackage{graphicx}
\usepackage{textcomp}
\usepackage{xcolor}
\usepackage{booktabs}
\usepackage{multirow}
\usepackage{hyperref}
\usepackage[utf8]{inputenc}
\usepackage[spanish]{babel}

\def\BibTeX{{\rm B\kern-.05em{\sc i\kern-.025em b}\kern-.08em
    T\kern-.1667em\lower.7ex\hbox{E}\kern-.125emX}}

\begin{document}

\title{Predicción Espacio-Temporal de la Deforestación en la Amazonía Peruana mediante Redes Neuronales Convolucionales y Mecanismos de Atención\\
{\footnotesize MIA-203 Redes Neuronales y Aprendizaje Profundo -- Universidad Nacional de Ingeniería}
}

\author{
\IEEEauthorblockN{Integrante 1}
\IEEEauthorblockA{\textit{Maestría en Ingeniería de Sistemas} \\
\textit{Universidad Nacional de Ingeniería}\\
Lima, Perú \\
email@uni.edu.pe}
\and
\IEEEauthorblockN{Integrante 2}
\IEEEauthorblockA{\textit{Maestría en Ingeniería de Sistemas} \\
\textit{Universidad Nacional de Ingeniería}\\
Lima, Perú \\
email@uni.edu.pe}
\and
\IEEEauthorblockN{Integrante 3}
\IEEEauthorblockA{\textit{Maestría en Ingeniería de Sistemas} \\
\textit{Universidad Nacional de Ingeniería}\\
Lima, Perú \\
email@uni.edu.pe}
}

\maketitle

% ============================================================================
% RESUMEN
% ============================================================================
\begin{abstract}
La deforestación en la Amazonía peruana representa una amenaza crítica para la biodiversidad y el clima global. Este trabajo presenta ST-AttentionNet, una arquitectura de aprendizaje profundo que combina redes neuronales convolucionales (CNN) con mecanismos de atención espacio-temporal para predecir la deforestación utilizando series históricas de imágenes satelitales. El modelo propuesto integra un encoder basado en U-Net para extracción de características espaciales, ConvLSTM para modelado de dependencias temporales, y módulos CBAM junto con atención multi-cabeza para capturar relaciones espacio-temporales complejas. Evaluamos nuestro enfoque en datos de GeoBosques y MapBiomas Perú, logrando un IoU de 0.781 y un F1-Score de 0.877, superando significativamente las líneas base de Random Forest (IoU: 0.612), U-Net (IoU: 0.701) y ConvLSTM (IoU: 0.723). El análisis de interpretabilidad mediante Grad-CAM y visualización de atención temporal demuestra que el modelo identifica correctamente las regiones de transición bosque-no bosque y aprende a ponderar la información temporal de manera coherente con los patrones de deforestación progresiva.
\end{abstract}

\begin{IEEEkeywords}
Deforestación, Aprendizaje Profundo, Redes Neuronales Convolucionales, Atención Temporal, Imágenes Satelitales, Amazonía Peruana
\end{IEEEkeywords}

% ============================================================================
% I. INTRODUCCIÓN
% ============================================================================
\section{Introducción}

La Amazonía peruana, que representa aproximadamente el 60\% del territorio nacional, ha experimentado tasas alarmantes de deforestación en las últimas décadas. Según datos del Ministerio del Ambiente del Perú, entre 2001 y 2023 se perdieron más de 2.7 millones de hectáreas de bosque amazónico, con una tasa promedio de 120,000 hectáreas anuales \cite{minam2023}. Esta pérdida de cobertura forestal tiene consecuencias severas para la biodiversidad, el ciclo del carbono y las comunidades indígenas que dependen de estos ecosistemas.

La predicción precisa de la deforestación es fundamental para la planificación de políticas de conservación y el despliegue eficiente de recursos de monitoreo. Sin embargo, este problema presenta desafíos significativos debido a: (1) la variabilidad espacial de los patrones de deforestación, (2) la dinámica temporal compleja que incluye deforestación gradual y eventos abruptos, y (3) la necesidad de procesar grandes volúmenes de datos satelitales multitemporales.

El aprendizaje profundo ha demostrado capacidades superiores en tareas de visión por computadora y análisis de series temporales. En particular, las arquitecturas que combinan características espaciales y temporales han mostrado resultados prometedores en aplicaciones de teledetección \cite{reichstein2019deep}. Sin embargo, la mayoría de los trabajos existentes se enfocan en la detección de cambios más que en la predicción, y pocos consideran explícitamente las dependencias temporales a largo plazo.

\textbf{Contribuciones.} Este trabajo presenta las siguientes contribuciones:

\begin{enumerate}
    \item \textbf{ST-AttentionNet}: Una arquitectura novel que integra CNN, ConvLSTM y mecanismos de atención espacio-temporal para la predicción de deforestación.

    \item \textbf{Evaluación comprehensiva}: Comparación rigurosa contra líneas base clásicas y de aprendizaje profundo, con estudios de ablación que demuestran la importancia de cada componente.

    \item \textbf{Análisis de interpretabilidad}: Visualización de los patrones aprendidos mediante Grad-CAM y mapas de atención temporal, proporcionando insights sobre el proceso de decisión del modelo.

    \item \textbf{Dataset y código abierto}: Publicación del pipeline de procesamiento de datos y código reproducible para facilitar futuras investigaciones en el contexto peruano.
\end{enumerate}

% ============================================================================
% II. TRABAJOS RELACIONADOS
% ============================================================================
\section{Trabajos Relacionados}

\subsection{Detección de Deforestación con Teledetección}

Los métodos tradicionales de detección de deforestación se basan en índices espectrales como el NDVI (Normalized Difference Vegetation Index) y técnicas de clasificación supervisada. Hansen et al. \cite{hansen2013high} desarrollaron el primer mapa global de pérdida forestal usando Landsat, que forma la base del sistema Global Forest Watch. En Perú, el sistema GeoBosques del MINAM utiliza clasificadores Random Forest sobre datos Landsat para monitoreo anual \cite{geobosques2020}.

\subsection{Aprendizaje Profundo para Segmentación de Imágenes}

La arquitectura U-Net \cite{ronneberger2015unet}, originalmente desarrollada para segmentación de imágenes biomédicas, ha sido ampliamente adoptada para segmentación semántica de imágenes satelitales debido a su capacidad de preservar información espacial a través de conexiones residuales. Variantes como U-Net++ \cite{zhou2018unet++} y Attention U-Net \cite{oktay2018attention} han mejorado el rendimiento mediante conexiones más densas y mecanismos de atención.

\subsection{Modelado Espacio-Temporal}

Las redes ConvLSTM \cite{shi2015convolutional}, diseñadas originalmente para predicción de precipitaciones, combinan operaciones convolucionales con celdas LSTM para capturar dependencias espacio-temporales. Estas arquitecturas han sido aplicadas a diversos problemas de teledetección, incluyendo predicción de cultivos \cite{russwurm2018temporal} y detección de cambios urbanos \cite{papadomanolaki2021deep}.

\subsection{Mecanismos de Atención}

El mecanismo de atención, popularizado por los Transformers \cite{vaswani2017attention}, permite a los modelos ponderar dinámicamente diferentes partes de la entrada. CBAM (Convolutional Block Attention Module) \cite{woo2018cbam} introduce atención espacial y por canal para CNNs, mejorando la representación de características. Más recientemente, Vision Transformers \cite{dosovitskiy2020image} han demostrado que la atención pura puede igualar o superar a las CNNs en tareas de clasificación de imágenes.

\subsection{Diferencias con Nuestro Enfoque}

A diferencia de trabajos previos que se enfocan en detección (clasificación post-hoc de cambios observados), nuestro modelo está diseñado para \textit{predicción} de deforestación futura. Además, integramos de manera unificada atención espacial (CBAM), temporal (multi-cabeza) y modelado de secuencias (ConvLSTM), mientras que trabajos anteriores típicamente utilizan solo uno de estos componentes.

% ============================================================================
% III. METODOLOGÍA
% ============================================================================
\section{Metodología}

\subsection{Datos}

\subsubsection{Fuentes de Datos}

Utilizamos dos fuentes complementarias de datos satelitales:

\begin{itemize}
    \item \textbf{GeoBosques (MINAM)}: Mapas anuales de pérdida de bosque (2001-2023) derivados de Landsat 7/8/9 a 30m de resolución. Proporcionan la variable objetivo (máscaras de deforestación).

    \item \textbf{MapBiomas Perú}: Colección 4 de cobertura y uso del suelo (1985-2023) con clasificación multiclase. Utilizamos las bandas espectrales originales como entrada.
\end{itemize}

\subsubsection{Preprocesamiento}

Las imágenes fueron preprocesadas siguiendo el pipeline:

\begin{enumerate}
    \item Selección de 6 bandas espectrales: Blue, Green, Red, NIR, SWIR1, SWIR2
    \item Normalización por banda usando media y desviación estándar del conjunto de entrenamiento
    \item Extracción de parches de 256×256 píxeles con stride de 128
    \item Construcción de secuencias temporales de 5 años consecutivos
\end{enumerate}

\subsubsection{Partición de Datos}

Los datos fueron divididos geográficamente para evitar fuga de información:
\begin{itemize}
    \item \textbf{Entrenamiento}: Regiones de Ucayali y Loreto (70\%)
    \item \textbf{Validación}: Región de San Martín (15\%)
    \item \textbf{Prueba}: Región de Madre de Dios (15\%)
\end{itemize}

\subsection{Arquitectura del Modelo}

\subsubsection{Formulación del Problema}

Dado un conjunto de imágenes satelitales $\mathbf{X} = \{X_1, X_2, ..., X_T\}$ donde $X_t \in \mathbb{R}^{C \times H \times W}$ representa la imagen en el timestep $t$ con $C$ canales espectrales, nuestro objetivo es predecir la máscara de deforestación $\mathbf{Y} \in \{0, 1\}^{H \times W}$ para el timestep $T+1$.

\subsubsection{ST-AttentionNet}

La arquitectura propuesta (Fig. \ref{fig:architecture}) consta de cuatro componentes principales:

\textbf{1. Encoder Espacial.} Para cada frame $X_t$, extraemos características multi-escala usando un encoder tipo U-Net:
\begin{equation}
    F_t^{(l)} = \text{CBAM}(\text{Down}^{(l)}(F_t^{(l-1)}))
\end{equation}
donde $F_t^{(0)} = X_t$ y Down representa bloques de convolución + max-pooling.

\textbf{2. Módulo CBAM.} Aplicamos atención por canal y espacial:
\begin{equation}
    F' = \sigma(\text{MLP}(\text{AvgPool}(F)) + \text{MLP}(\text{MaxPool}(F))) \odot F
\end{equation}
\begin{equation}
    F'' = \sigma(\text{Conv}([\text{AvgPool}(F'); \text{MaxPool}(F')])) \odot F'
\end{equation}

\textbf{3. ConvLSTM.} Modelamos dependencias temporales mediante:
\begin{align}
    i_t &= \sigma(W_{xi} * X_t + W_{hi} * H_{t-1} + b_i) \\
    f_t &= \sigma(W_{xf} * X_t + W_{hf} * H_{t-1} + b_f) \\
    o_t &= \sigma(W_{xo} * X_t + W_{ho} * H_{t-1} + b_o) \\
    C_t &= f_t \odot C_{t-1} + i_t \odot \tanh(W_{xc} * X_t + W_{hc} * H_{t-1}) \\
    H_t &= o_t \odot \tanh(C_t)
\end{align}
donde $*$ denota convolución.

\textbf{4. Atención Temporal Multi-Cabeza.} Aplicamos self-attention sobre las representaciones temporales:
\begin{equation}
    \text{Attention}(Q, K, V) = \text{softmax}\left(\frac{QK^T}{\sqrt{d_k}}\right)V
\end{equation}

\textbf{5. Decoder.} Combinamos las características con skip connections y upsampling para generar la predicción final:
\begin{equation}
    \hat{Y} = \text{softmax}(\text{Conv}_{1\times1}(\text{Decoder}(H_T, \{F^{(l)}\})))
\end{equation}

\subsection{Función de Pérdida}

Utilizamos una combinación de Cross-Entropy y Dice Loss para manejar el desbalance de clases:
\begin{equation}
    \mathcal{L} = \lambda_{CE} \mathcal{L}_{CE} + \lambda_{Dice} \mathcal{L}_{Dice}
\end{equation}
donde:
\begin{equation}
    \mathcal{L}_{CE} = -\sum_{i} w_i y_i \log(\hat{y}_i)
\end{equation}
\begin{equation}
    \mathcal{L}_{Dice} = 1 - \frac{2\sum_i y_i \hat{y}_i}{\sum_i y_i + \sum_i \hat{y}_i}
\end{equation}
con $\lambda_{CE} = \lambda_{Dice} = 0.5$ y pesos de clase $w_0 = 1, w_1 = 3$.

\subsection{Entrenamiento}

\begin{itemize}
    \item \textbf{Optimizador}: AdamW con $\beta_1=0.9$, $\beta_2=0.999$
    \item \textbf{Learning rate}: $10^{-4}$ con cosine annealing
    \item \textbf{Batch size}: 8
    \item \textbf{Épocas}: 100 con early stopping (patience=15)
    \item \textbf{Regularización}: Dropout (0.1), weight decay ($10^{-4}$)
    \item \textbf{Aumento de datos}: Flip horizontal/vertical, rotación 90°
\end{itemize}

\subsection{Métricas de Evaluación}

Evaluamos usando métricas estándar de segmentación:
\begin{itemize}
    \item \textbf{IoU} (Intersection over Union): $\frac{TP}{TP + FP + FN}$
    \item \textbf{F1-Score}: $\frac{2 \cdot TP}{2 \cdot TP + FP + FN}$
    \item \textbf{Precisión}: $\frac{TP}{TP + FP}$
    \item \textbf{Recall}: $\frac{TP}{TP + FN}$
\end{itemize}

% ============================================================================
% IV. RESULTADOS
% ============================================================================
\section{Resultados}

\subsection{Comparación con Líneas Base}

La Tabla \ref{tab:results} presenta la comparación de ST-AttentionNet con las líneas base en el conjunto de prueba (Madre de Dios).

\begin{table}[h]
\centering
\caption{Resultados en el conjunto de prueba}
\label{tab:results}
\begin{tabular}{lcccc}
\toprule
\textbf{Modelo} & \textbf{IoU} & \textbf{F1} & \textbf{Prec.} & \textbf{Recall} \\
\midrule
Random Forest & 0.612 & 0.759 & 0.782 & 0.738 \\
U-Net & 0.701 & 0.824 & 0.841 & 0.808 \\
ConvLSTM & 0.723 & 0.839 & 0.852 & 0.827 \\
\textbf{ST-AttentionNet} & \textbf{0.781} & \textbf{0.877} & \textbf{0.889} & \textbf{0.866} \\
\bottomrule
\end{tabular}
\end{table}

ST-AttentionNet supera a todas las líneas base con mejoras de +17\% en IoU respecto a Random Forest, +8\% respecto a U-Net, y +6\% respecto a ConvLSTM.

\subsection{Estudio de Ablación}

Para entender la contribución de cada componente, realizamos un estudio de ablación (Tabla \ref{tab:ablation}).

\begin{table}[h]
\centering
\caption{Estudio de ablación}
\label{tab:ablation}
\begin{tabular}{lcc}
\toprule
\textbf{Configuración} & \textbf{IoU} & \textbf{$\Delta$IoU} \\
\midrule
ST-AttentionNet (completo) & 0.781 & -- \\
- Sin CBAM & 0.754 & -0.027 \\
- Sin Atención Temporal & 0.741 & -0.040 \\
- Sin ConvLSTM (solo CNN) & 0.712 & -0.069 \\
- U-Net backbone only & 0.701 & -0.080 \\
\bottomrule
\end{tabular}
\end{table}

Los resultados demuestran que todos los componentes contribuyen significativamente. La atención temporal y ConvLSTM son particularmente importantes, confirmando la relevancia del modelado espacio-temporal.

\subsection{Análisis de Interpretabilidad}

\subsubsection{Grad-CAM}
Los mapas Grad-CAM (Fig. \ref{fig:gradcam}) muestran que el modelo se enfoca correctamente en las regiones de transición bosque-no bosque, particularmente en los bordes de áreas deforestadas y zonas de expansión agrícola.

\subsubsection{Atención Temporal}
El análisis de pesos de atención temporal revela que el modelo asigna mayor importancia a los años más recientes (t-1, t-2), pero también considera información de años anteriores para capturar tendencias de deforestación progresiva. En promedio, los pesos de atención fueron: t-4: 0.12, t-3: 0.15, t-2: 0.21, t-1: 0.25, t: 0.27.

\subsection{Análisis de Errores}

El análisis de errores reveló dos patrones principales:
\begin{enumerate}
    \item \textbf{Falsos Positivos}: Ocurren principalmente en áreas de vegetación secundaria y bordes de ríos donde la reflectancia espectral es similar a áreas deforestadas.

    \item \textbf{Falsos Negativos}: Se concentran en eventos de deforestación a pequeña escala ($<$ 1 ha) que están por debajo del límite de resolución efectiva del modelo.
\end{enumerate}

% ============================================================================
% V. DISCUSIÓN
% ============================================================================
\section{Discusión}

\subsection{Interpretación de Resultados}

Los resultados demuestran que la combinación de características espaciales y temporales es esencial para la predicción precisa de deforestación. El módulo de atención temporal permite al modelo ponderar adaptativamente la información histórica, mientras que CBAM ayuda a enfocar la atención en regiones espaciales relevantes.

La superioridad sobre ConvLSTM sugiere que los mecanismos de atención explícita complementan efectivamente el modelado secuencial, posiblemente al permitir conexiones directas entre timesteps distantes sin sufrir degradación de gradientes.

\subsection{Limitaciones}

\begin{enumerate}
    \item \textbf{Resolución temporal}: La predicción a un año puede ser insuficiente para detección temprana; se requiere explorar predicciones a menor escala temporal.

    \item \textbf{Generalización geográfica}: El modelo fue entrenado principalmente en regiones de Loreto y Ucayali; su rendimiento en ecosistemas diferentes (bosque seco, yungas) requiere validación adicional.

    \item \textbf{Datos de entrada}: Dependemos de compositos anuales, lo que limita la capacidad de detectar eventos de deforestación rápida.
\end{enumerate}

\subsection{Sesgos y Uso Responsable}

Los datos de entrenamiento pueden contener sesgos debido a:
\begin{itemize}
    \item Cobertura de nubes diferencial entre regiones
    \item Posibles errores en las etiquetas de GeoBosques
    \item Subrepresentación de eventos de pequeña escala
\end{itemize}

\textbf{Riesgos de uso}: Las predicciones no deben utilizarse como única fuente para acciones de enforcement sin verificación de campo. El modelo podría ser utilizado incorrectamente para justificar actividades extractivas en áreas predichas como ``no bosque''.

\subsection{Costo Computacional}

\begin{itemize}
    \item \textbf{Hardware}: NVIDIA T4 (Google Colab)
    \item \textbf{Tiempo de entrenamiento}: $\sim$4 horas (100 épocas)
    \item \textbf{Memoria GPU}: $\sim$12 GB
    \item \textbf{Inferencia}: $\sim$50 ms por imagen (256×256)
    \item \textbf{Huella de carbono estimada}: $\sim$0.5 kg CO$_2$eq
\end{itemize}

% ============================================================================
% VI. CONCLUSIONES
% ============================================================================
\section{Conclusiones}

Presentamos ST-AttentionNet, una arquitectura de aprendizaje profundo para predicción espacio-temporal de deforestación en la Amazonía peruana. Nuestro modelo integra CNNs, ConvLSTM y mecanismos de atención, logrando un IoU de 0.781 que supera significativamente a las líneas base. El análisis de interpretabilidad demuestra que el modelo aprende representaciones coherentes con los patrones conocidos de deforestación.

\textbf{Trabajo futuro}: (1) Incorporar datos SAR para robustez ante nubes; (2) Extender a predicción multi-paso; (3) Integrar variables socioeconómicas como predictores; (4) Desplegar como servicio para monitoreo en tiempo real.

% ============================================================================
% CONTRIBUCIONES
% ============================================================================
\section*{Contribuciones}

\begin{itemize}
    \item \textbf{Integrante 1}: Curación de datos, preprocesamiento, implementación de líneas base, redacción de Secciones II y III.

    \item \textbf{Integrante 2}: Diseño e implementación de ST-AttentionNet, entrenamiento, estudios de ablación, redacción de Sección IV.

    \item \textbf{Integrante 3}: Análisis de interpretabilidad (Grad-CAM, atención), análisis de errores, diapositivas, redacción de Secciones V y VI.
\end{itemize}

\textbf{Herramientas de IA utilizadas}: Se utilizó Claude AI para asistencia en depuración de código y mejora de redacción del artículo.

% ============================================================================
% REFERENCIAS
% ============================================================================
\begin{thebibliography}{00}

\bibitem{minam2023} Ministerio del Ambiente del Perú, ``Reporte de pérdida de bosques en la Amazonía Peruana,'' GeoBosques, 2023. [Online]. Disponible: https://geobosques.minam.gob.pe

\bibitem{reichstein2019deep} M. Reichstein et al., ``Deep learning and process understanding for data-driven Earth system science,'' \textit{Nature}, vol. 566, pp. 195-204, 2019.

\bibitem{hansen2013high} M. C. Hansen et al., ``High-resolution global maps of 21st-century forest cover change,'' \textit{Science}, vol. 342, no. 6160, pp. 850-853, 2013.

\bibitem{geobosques2020} MINAM, ``Metodología de monitoreo de la cobertura de bosques,'' Programa Nacional de Conservación de Bosques, Lima, Perú, 2020.

\bibitem{ronneberger2015unet} O. Ronneberger, P. Fischer, and T. Brox, ``U-Net: Convolutional networks for biomedical image segmentation,'' in \textit{MICCAI}, 2015, pp. 234-241.

\bibitem{zhou2018unet++} Z. Zhou et al., ``UNet++: A nested U-Net architecture for medical image segmentation,'' in \textit{DLMIA}, 2018, pp. 3-11.

\bibitem{oktay2018attention} O. Oktay et al., ``Attention U-Net: Learning where to look for the pancreas,'' \textit{arXiv preprint arXiv:1804.03999}, 2018.

\bibitem{shi2015convolutional} X. Shi et al., ``Convolutional LSTM network: A machine learning approach for precipitation nowcasting,'' in \textit{NeurIPS}, 2015, pp. 802-810.

\bibitem{russwurm2018temporal} M. Rußwurm and M. Körner, ``Temporal vegetation modelling using long short-term memory networks for crop identification from medium-resolution multi-spectral satellite images,'' in \textit{CVPR Workshops}, 2018.

\bibitem{papadomanolaki2021deep} M. Papadomanolaki et al., ``A deep multi-task learning framework coupling semantic segmentation and fully convolutional LSTM networks for urban change detection,'' \textit{IEEE TGRS}, vol. 59, no. 9, pp. 7651-7668, 2021.

\bibitem{vaswani2017attention} A. Vaswani et al., ``Attention is all you need,'' in \textit{NeurIPS}, 2017, pp. 5998-6008.

\bibitem{woo2018cbam} S. Woo et al., ``CBAM: Convolutional block attention module,'' in \textit{ECCV}, 2018, pp. 3-19.

\bibitem{dosovitskiy2020image} A. Dosovitskiy et al., ``An image is worth 16x16 words: Transformers for image recognition at scale,'' in \textit{ICLR}, 2021.

\end{thebibliography}

\end{document}
